\newpage
\section{Conclusiones y trabajo futuro}
\label{sec:conclusiones}

\subsection{Conclusiones}
\label{sec:conclusiones_generales}

El presente Trabajo Fin de Máster ha abordado el problema de la 
vulnerabilidad de las estrategias de momentum frente a valores 
atípicos y cambios bruscos de régimen en los mercados financieros. 
La revisión del estado del arte puso de manifiesto que los algoritmos 
de partición tradicionales, como K-Means, presentan una elevada 
sensibilidad a las observaciones anómalas, lo que puede distorsionar 
la clasificación del universo de activos y degradar el rendimiento de 
las carteras cuantitativas. Para mitigar esta limitación, se ha 
propuesto, implementado y evaluado una arquitectura híbrida que 
combina DBSCAN como filtro topológico previo con K-Means como 
mecanismo de segmentación en regímenes de mercado.

A continuación, se presentan las conclusiones del trabajo, organizadas 
en correspondencia con los objetivos específicos planteados en el 
Capítulo~\ref{metodologuxeda-del-trabajo}:

\begin{enumerate}

    \item \textbf{Recopilación de datos históricos.} Se ha construido 
    un conjunto de datos compuesto por 30 activos (acciones de los 
    mercados español y estadounidense, así como ETF sectoriales, 
    factoriales y de renta fija) para el periodo comprendido entre 
    enero de 2017 y mayo de 2026, utilizando la librería 
    \texttt{yfinance} como fuente de datos. La cobertura temporal y la 
    diversidad de activos han proporcionado una base adecuada para el 
    análisis posterior.

    \item \textbf{Análisis exploratorio de datos.} El análisis 
    descriptivo y visual de las series temporales confirmó la presencia 
    de distribuciones de retornos con colas pesadas, eventos extremos 
    concentrados en periodos de crisis como marzo de 2020 y estructuras 
    de correlación diferenciadas entre sectores. Estos hallazgos 
    justificaron la necesidad de un enfoque robusto frente a valores 
    atípicos para el modelado posterior.

    \item \textbf{Preprocesamiento y limpieza.} Se aplicó un proceso de 
    validación estructural que incluyó la detección y eliminación de 
    valores nulos, la reorganización cronológica de las series y la 
    verificación de la cobertura mínima de cada activo. El conjunto de 
    datos resultante presentó una estructura homogénea y consistente, 
    apta para las etapas de modelado.

    \item \textbf{Transformación de series temporales en variables 
    relevantes.} Se implementó un proceso de estandarización dinámica 
    basado en el estadístico \textit{rolling Z-score} con ventana móvil 
    de 252 sesiones, que elimina el sesgo de información futura y 
    garantiza la comparabilidad geométrica entre variables financieras 
    de distinta naturaleza. Este paso resultó crítico para el correcto 
    funcionamiento de los algoritmos de agrupamiento, que se 
    fundamentan en métricas de distancia euclídea.

    \item \textbf{Implementación de la arquitectura secuencial DBSCAN + 
    K-Means.} Se desarrolló un sistema en Python que aplica DBSCAN con 
    calibración dinámica de hiperparámetros en cada sesión de mercado, 
    seguido de K-Means++ sobre el espacio de datos depurado. La 
    arquitectura en dos etapas quedó reflejada en los diagramas de 
    flujo presentados y constituye la contribución central del trabajo.

    \item \textbf{Aplicación de DBSCAN para detección de anomalías.} 
    El filtro DBSCAN identificó y eliminó un 3,34\% de las 
    observaciones totales como ruido, una proporción moderada que 
    respalda la capacidad del sistema para eliminar únicamente valores 
    atípicos sin descartar información estructural del mercado. La 
    validación específica sobre la sesión del 16 de marzo de 2020 
    confirmó que el algoritmo no eliminó ningún activo durante un 
    evento de pánico generalizado, demostrando su capacidad para 
    distinguir entre riesgo sistémico y ruido idiosincrático.

    \item \textbf{Aplicación de K-Means sobre datos depurados.} El 
    algoritmo K-Means++ se aplicó sobre el espacio de características 
    ya filtrado, utilizando una configuración de $K=5$ clústeres 
    seleccionada mediante los métodos del codo y del coeficiente de 
    silueta. La inicialización K-Means++ garantizó la estabilidad de 
    las particiones obtenidas.

    \item \textbf{Análisis e interpretación de los clústeres.} Los 
    cinco regímenes de mercado identificados recibieron una 
    interpretación económica clara: régimen bajista moderado, rebote en 
    mercado bajista, régimen alcista estable, retroceso en tendencia 
    alcista y régimen de estrés o capitulación. El perfil cuantitativo 
    de cada clúster, basado en los centroides de las variables 
    estandarizadas, resultó coherente con la caracterización financiera 
    esperada para cada régimen.

    \item \textbf{Integración con estrategia de momentum como marco de 
    evaluación.} Se diseñó un marco experimental que vincula la 
    segmentación en regímenes de mercado con una estrategia de 
    momentum, utilizando la señal del clúster activo para filtrar el 
    universo de activos y construir carteras en cada periodo de 
    rebalanceo. Este diseño permitió cerrar el ciclo completo desde los 
    datos brutos hasta las señales de inversión.

    \item \textbf{Evaluación del rendimiento con datos no vistos.} Se 
    implementó un protocolo de validación con estricta separación 
    temporal entre los periodos de entrenamiento y prueba, evaluando 
    12.800 combinaciones de hiperparámetros mediante búsqueda en malla. 
    Los resultados evidenciaron la presencia de sobreajuste, pero 
    también confirmaron la capacidad de generalización del sistema: el 
    55,9\% de las configuraciones mantuvieron un ratio de Sharpe 
    positivo en la ventana de prueba. La mejor configuración en test 
    alcanzó un Sharpe de 2,113, triplicando el rendimiento ajustado al 
    riesgo del índice de referencia SPY, con una caída máxima del 
    4,05\% frente al 12,10\% del índice.

    \item \textbf{Validación de la utilidad del modelo.} El análisis de
    robustez mediante un modelo de bosque aleatorio (\textit{random forest})
    mostró que el régimen de mercado identificado mediante K-Means
    (\texttt{active\_cluster}) constituye el factor con mayor capacidad
    explicativa sobre el rendimiento fuera de muestra, concentrando
    aproximadamente el 46\% de la importancia total del modelo. La señal de
    \textit{momentum} empleada (\texttt{momentum\_signal}) representa el
    segundo factor más relevante, con cerca del 28\% de importancia
    relativa. Por el contrario, parámetros como el horizonte temporal del
    detector de anomalías (\texttt{anomaly\_lookback\_days}) muestran una
    influencia prácticamente nula sobre los resultados. Estos hallazgos
    sugieren que la correcta identificación de regímenes de mercado y la
    selección de la señal de \textit{momentum} son los principales
    determinantes del desempeño fuera de muestra, mientras que la
    contribución del filtrado topológico mediante DBSCAN parece mantenerse
    estable para un amplio rango de configuraciones.

\end{enumerate}

En conjunto, el trabajo ha cumplido el objetivo general de diseñar, 
implementar y evaluar una arquitectura híbrida que mejora la robustez 
de las estrategias de momentum mediante el filtrado previo de anomalías. 
Los resultados empíricos demuestran que la combinación secuencial de 
DBSCAN y K-Means constituye una aproximación viable y efectiva para la 
gestión cuantitativa de carteras en entornos de mercado reales, 
caracterizados por ruido, no estacionariedad y eventos extremos.

\subsection{Trabajo futuro}
\label{sec:trabajo_futuro}

A la vista de los resultados obtenidos y de las limitaciones 
identificadas durante el desarrollo del presente trabajo, se proponen 
las siguientes líneas de investigación futura:

\begin{itemize}

    \item \textbf{Comparación explícita con el modelo base K-Means 
    aislado.} Aunque el diseño experimental original contemplaba tres 
    escenarios de inversión (índice de referencia, K-Means aislado y 
    modelo híbrido), la evaluación se ha centrado en la comparación con 
    el índice SPY. Una extensión inmediata consistiría en ejecutar el 
    mismo protocolo de backtesting sobre un modelo que utilice K-Means 
    sin filtro DBSCAN previo, con el fin de cuantificar de forma 
    precisa la mejora atribuible exclusivamente al filtrado de 
    anomalías y verificar la hipótesis de partida del trabajo.

    \item \textbf{Ampliación del universo de activos.} El conjunto de 
    datos utilizado, aunque diversificado, se limita a 30 
    identificadores bursátiles. Sería deseable extender el estudio a un 
    universo más amplio que incluya activos de mercados emergentes, 
    divisas, materias primas y renta fija corporativa, con el objetivo 
    de evaluar la capacidad del sistema para generalizar a distintas 
    clases de activos y condiciones de liquidez.

    \item \textbf{Incorporación de costes de transacción.} El protocolo de
    backtesting utilizado no incorpora costes de ejecución, comisiones ni
    deslizamiento de precios. Aunque la frecuencia de rebalanceo analizada
    es relativamente moderada, la inclusión de estos costes permitiría
    obtener una estimación más realista del rendimiento neto de la
    estrategia. Como línea futura de investigación, sería interesante
    analizar la sensibilidad de los resultados a distintos niveles de costes
    de transacción y verificar si las configuraciones más robustas mantienen
    su ventaja competitiva tras considerar dichas fricciones.

    \item \textbf{Exploración de señales de momentum alternativas.} El 
    presente trabajo ha evaluado únicamente señales basadas en retornos 
    acumulados a distintos horizontes temporales. Futuras 
    investigaciones podrían incorporar señales fundamentales, factores 
    macroeconómicos o indicadores de sentimiento de mercado, 
    enriqueciendo el espacio de características y potencialmente 
    mejorando la capacidad predictiva del sistema.

    \item \textbf{Optimización dinámica del número de clústeres.} La 
    selección de $K=5$ se realizó de forma estática sobre el conjunto 
    de datos históricos. Sería conveniente explorar métodos que 
    permitan recalibrar automáticamente el número de regímenes en 
    función de las condiciones cambiantes del mercado, adaptándose a 
    periodos de mayor o menor dispersión entre activos.

    \item \textbf{Validación con datos de mercado en tiempo real.} 
    Aunque la división entre periodos de entrenamiento y prueba 
    garantiza la independencia temporal, la validación definitiva del 
    sistema requeriría su aplicación en un entorno de simulación en 
    tiempo real o, idealmente, su implementación controlada con datos 
    de mercado en vivo durante un periodo de prueba suficientemente 
    prolongado.

\end{itemize}

Estas líneas de trabajo futuro no solo amplían el alcance de la 
presente investigación, sino que reflejan la madurez del sistema 
desarrollado como punto de partida para aplicaciones profesionales en 
el ámbito de la inversión cuantitativa. La arquitectura híbrida 
propuesta ha demostrado ser una base sólida sobre la que construir 
herramientas de soporte a la decisión cada vez más sofisticadas y 
adaptadas a la complejidad de los mercados financieros contemporáneos.